% ==============================================================
%  DADOS DO RELATÓRIO
%  Edite apenas este arquivo para configurar a capa e o cabeçalho.
% ==============================================================


% --- TÍTULO ---
% Aparece em destaque na capa.
% Para usar um título curto diferente no cabeçalho das páginas:
%   \titulo[Título Curto]{Título Completo da Capa}
% Se omitir o argumento opcional, o mesmo texto é usado nos dois lugares.

\titulo{Relatório de Sistemas Digitais}


% --- SUBTÍTULO (opcional) ---
% Aparece abaixo do título na capa, em destaque menor.
% Comente ou remova esta linha para omitir o subtítulo.

\subtitulo{Simulação do Processador LEGv8}


% --- AUTORES ---
% Um autor por linha, separados por \\.
% Formato sugerido:  Nome Completo -- NUSP 1234567

\Autores{
    Nome do Aluno -- NUSP XXXXXXX\\
    Nome do Aluno -- NUSP XXXXXXX\\
    Nome do Aluno -- NUSP XXXXXXX\\
    Nome do Aluno -- NUSP XXXXXXX
}


% --- DEPARTAMENTO ---
% Aparece no cabeçalho da capa, abaixo do logo.

\departamento{Departamento de Computação e Sistemas Digitais}


% --- DISCIPLINA (opcional) ---
% Aparece na capa, logo acima da data.
% Descomente a linha abaixo para ativar.

% \disciplina{PCS3635 -- Laboratório de Sistemas Digitais I}


% --- DATA ---
% \today insere a data de compilação automaticamente.
% Para uma data fixa escreva, por exemplo:  \data{maio de 2026}

\data{\today}
